\documentclass[11pt,a4paper]{article}
\usepackage[margin=25mm]{geometry}
\usepackage{amsmath,amssymb,amsthm}
\newtheorem{theorem}{Theorem}
\title{Comparing two complete lists of marked groups}
\author{Human Compression project: L071 recovery}
\date{Working proof; distinct review pending}
\begin{document}\maketitle
Let $G\trianglelefteq A\le\operatorname{Sym}(\Omega)$, and suppose
conjugation induces every automorphism of $G$. The kernel of this action
is allowed to be nontrivial. Let $\mathcal T$ be a finite list containing
one representative of every automorphism orbit of the desired ordered
generating tuples of $G$. Its completeness is a premise supplied by the
acquired enumeration theorem and its completed execution.

Store $\mathcal T$ in a rooted tree whose edges carry permutations.
Each path from the root records a literal initial segment of a tuple;
identical initial segments have the same node. A terminal node records
the tuple's rank and its position in the enumeration. When the list comes
from the acquired stabilizer recursion, distinct children at a node lie
in distinct orbits of the pointwise stabilizer of the preceding tuple.
The comparison below also verifies each final match explicitly.

Let $H$ be a second concrete group, let $f:H\longrightarrow G$ be a
verified isomorphism, and let $\mathcal S$ be a finite list of tuples
in $H$. The two groups need not have the same permutation degree.

\begin{theorem}
Suppose each $s\in\mathcal S$ is mapped, by $f$ followed by a recorded
conjugation in $A$, to a terminal tuple of $\mathcal T$. Suppose these
terminal tuples are distinct, and $|\mathcal S|=|\mathcal T|$.
Then $\mathcal S$ and $\mathcal T$ give precisely the same classes
of ordered marked groups, each once. All properties of the tuples that
are invariant under group isomorphism transfer between the two lists.
\end{theorem}
\begin{proof}
Conjugation by an element of $A$ restricts to an automorphism of $G$.
Each recorded equality therefore gives an isomorphism from the source
marked group to its named terminal tuple. Distinct terminal tuples
represent distinct automorphism orbits by the acquired enumeration
theorem. Thus the resulting assignment is injective on marked classes.
The equality of the two finite list sizes makes it surjective. In
particular it transfers generation, the orders of adjacent products,
the intersection property and any stated equivalence with a reversed
marking. The assertion concerns ordered isomorphism classes; no reversal
is silently added to their equivalence relation.
\end{proof}

Here is an explicit way to find and certify the conjugations. Start with
$t=f(s)$, $a=1$ and $C=A$. At a node with children marked by $r$, find
$h\in C$ such that $t_i^h=r$ for the next position $i$. Replace the
whole tuple by $t^h$, replace $a$ by $ah$, and replace $C$ by
$C_C(r)$. Move to the child marked $r$ and continue. If no such child
exists, the comparison fails; it does not discard the source tuple.

Inductively, $C$ centralizes every already matched mark. Thus the new
conjugation preserves the entire matched initial segment, and the updated
centralizer is exactly its pointwise stabilizer after adjoining $r$.
For right conjugation, the accumulated element is $a=h_1\cdots h_i$.
At the end the checker verifies both $a\in A$ and the complete literal
equality $f(s)^a=t_{\mathrm{terminal}}$. It also checks that the named
terminal node exists and has not occurred for an earlier source tuple.
These checks are the hypotheses of the theorem; a successful orbit-search
call alone is not substituted for the final equality.

\paragraph{The concrete source comparison.}
The auxiliary archive supplied by Hartley and Hulpke records permutation
generators and lists of indices selecting each marking. Our reader accepts
only the two empty-list initializations, literal permutation additions,
and tuple additions using already defined indices. It reconstructs data
from these expressions and does not execute the source's loader program.
The original archive and every parsed file are retained with their exact
hashes. They are comparison inputs acquired after our enumeration; neither
their counts nor their generators direct its search.

For a nonempty source list, take $H$ to be generated by its first tuple.
The execution checks $|H|=|G|$ and membership in $H$ of every source
generator, constructs $f$ by an exact group-isomorphism algorithm, and
checks that $f$ is bijective. It records generators of $H$, their images,
all source-generator images and the individual conjugations. Empty source
lists are compared only with empty completed enumerations. Thus absence
claims retain the same completeness premise as positive classifications.

\paragraph{Verification boundary and attribution.}
The finite group computations use the retained exact GAP algorithms for
group isomorphism, membership, centralizers and conjugating elements.
They are computational premises, not Lean declarations. This comparison
does not establish the completeness of $\mathcal T$ from agreement of
totals; its complete enumeration proof remains a separate paid premise.
The source markings and their classification are attributed to Hartley
and Hulpke, and the subsequent verification to Leemans and Mixer.
The present contribution is a recorded bijection from an independent
enumeration to those original markings, using the same stabilizer structure
as the common search theorem. No new classification priority is claimed.
\end{document}
